% =============================================================================
% Appendix: knowledge-area annotation (simplified version).
%
% Condensed version of appendix_knowledge_area.tex; labels are distinct, so both
% files can be included in the same document. Usage: \input after \appendix.
% Required packages: graphicx, booktabs, natbib (\citep), hyperref or url,
% and xcolor (as used in the manuscript). Bibliography entries
% are in ablations/knowledge_area/references_knowledge_area.bib.
%
% Every number comes from the run artifacts under
% output/knowledge-annotation-work/ (full-run, full-run-2,
% uncertain-taxonomy-v1.3, uncertain-taxonomy-v1.4, full-run-taxonomy-v1.4 and
% its post-annotation/ directory), and the published test/dev snapshots under
% output/fewshot-work/main/datasets. These populations are reported separately by
% ablations/knowledge_area/make_figures.py, which also generates the figures;
% the sizes of the audits behind each taxonomy version come from
% src/mmlu_pt/annotation/knowledge_area/TAXONOMY_CHANGELOG.md. By default
% figures are read from ablations/knowledge_area/figures; to change this, define the macro before the
% \input, e.g. \newcommand{\KnowledgeAreaFigDir}{ablations/knowledge_area/figures}
% =============================================================================
\providecommand{\KnowledgeAreaFigDir}{ablations/knowledge_area/figures}

\section{Knowledge-Area Annotation}
\label{app:ka-short}

Every released question carries a \emph{subject}, chosen from a closed list of
candidates fixed per exam, and a \emph{macro area}, derived from the subject
by a table and never generated. The paper uses them to report accuracy per
area and to stratify the few-shot demonstrations by macro area and academic
level. The labels are assigned by an open-weight language model that reads
the question, its choices, the exam, and the edition, chooses one candidate or
abstains with \texttt{UNCERTAIN}, and does so twice with independent seeds. A
third request by the same model settles the disagreements. The method was
fixed before the first full run and never changed, what changed was the
taxonomy. Five versions were written (1.0 to 1.4), three were run over the
annotation corpus (1.1, 1.2, and 1.4), and two over the abstentions of the previous
full run (1.3 and 1.4), as Table~\ref{tab:ka-short-overview} summarizes.

\begin{table}[t]
\centering
\footnotesize
\setlength{\tabcolsep}{3.5pt}
\resizebox{\linewidth}{!}{
\begin{tabular}{@{}lrrrrlrrrr@{}}
\toprule
Ver. & Subj. & Areas & Exams & Rules & Run & $n$ & Agr. & Adj. & Unc. \\
\midrule
1.0 & 70 & 11 & 15 & 8 & -- & -- & -- & -- & -- \\
1.1 & 70 & 11 & 17 & 8 & full & 41,653 & 91.6 & 8.4 & 88 \\
1.2 & 78 & 9 & 17 & 14 & full & 41,653 & 92.1 & 7.9 & 23 \\
1.3 & 78 & 9 & 17 & 15 & subset & 23 & -- & -- & 9 \\
1.4 & 78 & 9 & 17 & 15 & subset & 15 & -- & -- & 1 \\
1.4 & 78 & 9 & 17 & 15 & full & 41,636 & 92.1 & 7.9 & 8 \\
\bottomrule
\end{tabular}
}
\caption{Taxonomy versions and the runs made with each. \emph{Subj.}, \emph{Areas}, \emph{Exams}, \emph{Rules}: subjects, macro areas, exams with a candidate list and annotation rules in the version; \emph{Run}: a full run over the corpus or a run over the abstentions of the previous full run; $n$: questions; \emph{Agr.}: percentage of questions on which the two independent passes chose the same subject; \emph{Adj.}: percentage sent to adjudication; \emph{Unc.}: questions left \texttt{UNCERTAIN}. Runs 1.1 and 1.2 annotated the filtered corpus; run 1.4 annotated the 41,636 questions left after the dataset revision. Agreement is the self-consistency of one model, not accuracy.}
\label{tab:ka-short-overview}
\end{table}

\paragraph{Summary.} (1)~One open-weight model, Qwen3.5-122B-A10B in FP8,
labels each question twice with independent seeds under a JSON-schema grammar
that admits only the exam's candidates and \texttt{UNCERTAIN}. A third request
by the same model adjudicates when the passes differ or either abstains.
(2)~The taxonomy grew from 70 subjects in 11 macro areas and 15 exams (1.0) to
78 subjects in 9 areas and 17 exams with 15 annotation rules (1.4), each step
motivated by an audit of the previous run's abstentions and low-confidence
labels. (3)~Abstentions fell from 88 (1.1) to 23 (1.2) to 8 (1.4) and to 0
after six manual assignments and two removals. (4)~Between consecutive full
runs, 8.0\% and 4.5\% of the labels changed, against a disagreement of about
8\% between the two passes of a single run, 549 questions took one of the
eight subjects added in 1.2. (5)~All nine macro areas of the test set
have at least 2,224 questions, and all 78 subjects are used, 15 of them with
fewer than 100 questions. (6)~The final run took 3\,h\,39\,min on one GPU and
processed 311.2 million prompt tokens.

% -----------------------------------------------------------------------------
\subsection{Annotation method}
\label{app:ka-short-method}

\paragraph{Model choice.} The annotator is Qwen3.5-122B-A10B-FP8, a
mixture-of-experts model with 122 billion parameters of which 10 billion are
active per token, in the FP8 weights published by its authors (revision
\texttt{a099dee7}\ldots), with its reasoning mode disabled.\footnote{\url{https://huggingface.co/Qwen/Qwen3.5-122B-A10B-FP8}, revision \texttt{a099dee70ccfcd8d5dda56aaa0b60cb8ecadabc9}, Apache 2.0 license. The model card is the source of the parameter counts, the language coverage, and the sampling guidance quoted here.} Four considerations led to it. (i)~Large instruction-tuned
models are an established substitute for crowd and expert annotators on
closed-label text classification, including zero-shot classification from a
list of labels with definitions
\citep{ding2023gpt3,he2024annollm,tan2024survey},
with the caveat that their labels must be validated per task
\citep{gligoric2025unconfident}, this appendix therefore reports consistency
rather than accuracy and the stage ends with a manual review. (ii)~The
annotator must not be one of the systems the benchmark evaluates: language
models used as judges favour their own outputs, and the bias grows with their
ability to recognize them \citep{panickssery2024llm,liu2024narcissistic}. All
19 open-weight checkpoints in the evaluation matrix have at most 35 billion
parameters and the 122B-A10B model is not among them, although six of them
belong to the same Qwen3.5 family. (iii)~Open weights with a pinned revision
make the annotation reproducible and auditable, which an API model that can
change silently does not offer, and open models have been shown to be viable
annotators \citep{alizadeh2025opensource}. The Qwen family documents broad
multilingual coverage, 119 languages including Portuguese for Qwen3
and 201 languages and dialects for Qwen3.5 according to their technical
report and model card, and is the base of the Portuguese Tucano~2 models
evaluated in the paper. (iv)~The model fits the
available hardware: the FP8 weights occupy 115.37~GiB of the 191.5~GB of one
NVIDIA B200, leaving a KV cache of 1,413,120 tokens for the 50 concurrent
sequences of 20,480 tokens used by the run, so a single GPU annotates the
corpus in under four hours while the other GPUs of the evaluation budget stay
free. The sampling parameters are the ones the model card recommends for its
thinking mode (temperature~1.0, top-$p$~0.95, top-$k$~20, presence
penalty~1.5), they were kept when the reasoning mode was switched off.

\paragraph{Prompt and decoding.} The system message states the task
(identify the one academic discipline whose knowledge is primarily required
to determine the correct answer), forbids solving the question or classifying
by superficial words or setting, states that exam and edition are contextual
priors and not automatic labels, allows \texttt{UNCERTAIN} when no candidate
fits, declares the question and the choices untrusted data, and asks for a
JSON object with a subject, an alternative subject, a confidence in $[0,1]$
and a justification of at most 400 characters. The user message gives the
exam and the edition, the annotation rules of the taxonomy version, the
numbered candidates of the exam with their definitions, and the question and
its choices as JSON, the answer key is never shown. The adjudication prompt
appends the two proposals, reduced to subject and justification, and asks the
model to assess them independently, choosing either, another candidate or
\texttt{UNCERTAIN}. Inference runs offline in vLLM~0.25.0 \citep{kwon2023efficient} under a JSON schema
whose subject fields are enumerations over the candidates and
\texttt{UNCERTAIN}, enforced by XGrammar \citep{dong2024xgrammar}. The
parsed object is validated locally and a failed request is retried with a new
seed, at most five times, with every seed derived from the base seed, the
question, the stage, the pass and the attempt. In the final run 86,571
requests succeeded, six were retried after a schema violation and all ended
at a stop token. Every attempt and token count is written to a SQLite
checkpoint, and a run resumes only under the same source, model revision,
taxonomy, prompt and code hashes.

\paragraph{Two passes and adjudication.} The labels of a language model
vary with the prompt and the model, and identical inputs can receive
different labels under sampling, a variance that a single sample per item
hides \citep{song2025good}, so every
question is labelled twice with independent seeds, in the spirit of
self-consistency checks that sample a model several times
\citep{wang2023selfconsistency}. When the two subjects agree and neither is
\texttt{UNCERTAIN}, the label is accepted with the smaller of the two
confidences, otherwise the adjudication request decides (status
\emph{adjudicated}, or \emph{uncertain} when it abstains). In the final run
the passes agreed on 92.1\% of the questions and 3,299 (7.9\%) went to
adjudication, where the model sided with pass~1 in 1,574 cases, with pass~2 in
1,640, chose a third subject in 77 and abstained in 8. Agreement between
passes is the self-consistency of one model and does not measure accuracy.

% -----------------------------------------------------------------------------
\subsection{Taxonomy versions and their effect}
\label{app:ka-short-taxonomy}

Four properties hold in every version. The macro area is a function of the
subject and is never offered to the model. Candidate lists only restrict: an
exam's list is a subset of the global list, no rule may depend on the answer,
and ENADE, whose booklets cover every degree programme, admits all subjects.
Each subject has a definition that names the neighbouring subjects it
excludes (Portuguese Language against Literature, Logic against Mathematics
and Programming, Internal Medicine against General Surgery, the branches of
Law against each other). And \texttt{UNCERTAIN} is a first-class label, so
that a question whose discipline is missing from the list or whose statement
is unreadable is not forced into a neighbour. Law is a macro area of its own
with 14 subjects, because the bar, magistracy and civil-service exams make it
the largest field of the corpus. Logic is separate from Mathematics and
Programming because the informatics olympiad consists largely of puzzles
solved from rules given in the statement.

Version 1.0 covered the 15 exams of the corpus at the time. 1.1 adds
candidate lists for ENEM (15 subjects) and IME (5) and two candidates to
existing lists, and is the baseline run. The audit of run 1.1 read its 88
abstentions and 2,901 labels with confidence below 0.90 and found disciplines
missing from the taxonomy, disciplines present but absent from an exam's
list, and ambiguous definitions and rules. Version 1.2 adds eight subjects
(among them radiology, occupational therapy, gastronomy and theology), merges
11 macro areas into 9, widens the lists of BACEN (18 to 26 candidates), BNDES
(30 to 36), the residency exams (5 to 11), POSCOMP (9 to 12), BLUEX (10 to
11) and ENADE (70 to 78), refines 29 definitions and raises the rules from 8
to 14. The analysis of run 1.2 read its 23 abstentions: eight were
unrecoverable statements, later removed by the dataset revision, and the rest
pointed to three candidate gaps (English comprehension in CNU, professional
ethics in the medical exams, social-security law in the residency exams),
nine definitions and a rule requiring the chosen label to match the
discipline named in the justification. Version 1.3 made those changes and, run on the 23 abstentions,
labelled 14. Version 1.4 refines the definition of English Language and three
rules, labelled 14 of the 15 questions left after the revision, and was frozen
before the third full run, with its hash recorded in the run manifest.

\paragraph{Effect on the labels.} Of the 88 abstentions of run 1.1, 78
received a subject in run 1.2. The 8 abstentions of run 1.4 are questions no
earlier run had abstained on. On the 41,636 questions common to the three
full runs, 3,316 labels (8.0\%) changed from 1.1 to 1.2 and 1,855 (4.5\%)
from 1.2 to 1.4, against a disagreement between the two passes of a single
run of 7.8\% (1.2) and 7.9\% (1.4). The change from 1.1 to 1.2 exceeds that
level only in the exams whose candidate lists 1.2 widened (BACEN 21.6\%, the
residency exams 15.7\%, POSCOMP 9.9\%), and 549 questions carry one of the
eight new subjects. The macro area is more stable than the subject: 96.4\% of
the questions keep it from 1.1 to 1.2 and 97.8\% from 1.2 to 1.4. Whether the
new labels are right is a separate question that these numbers do not answer:
the audits behind the versions are engineering reviews of the model's output,
not expert annotation.

% -----------------------------------------------------------------------------
\subsection{Final test set}
\label{app:ka-short-dataset}

\paragraph{Population.} The published benchmark contains 37,800 test
questions, 8,705 at high-school and 29,095 at undergraduate level. Its 75 dev
examples (30 high-school and 45 undergraduate) provide the few-shot
demonstrations and are excluded from all test distributions below.
The test questions are selected from the revised deduplicated corpus and
carry the labels of annotation run~1.4, joined by their content identifiers
(Appendix~\ref{app:rev-short}).

\paragraph{By level and macro area.} The two levels have different profiles.
High-school questions are 32.3\% Mathematics, Statistics and Logic
(driven by ENEM and OBI), 26.7\% Languages, 23.1\% Natural Sciences and
16.8\% Humanities, with 80 programming items from OBI and 17 Physical
Education items as the only labels outside those areas. Law, Economics and
Social Sciences have no high-school question, since no vestibular list
contains their subjects. Undergraduate questions are 36.0\% Law,
followed by Health Sciences (13.7\%), Economics, Business and Accounting
(11.5\%) and Computing, Engineering, Architecture and Design (11.3\%). Over
the whole test set Law holds 27.7\% of the questions and the smallest area,
Social and Applied Social Sciences, 5.9\% (2,224 questions).

\paragraph{By subject.} All 78 subjects are used (Table~\ref{tab:ka-short-subjects}).
Mathematics is the largest with 2,561 questions, followed by Civil Law and
Civil Procedure (2,317), Criminal Law and Criminal Procedure (1,749),
Portuguese Language (1,623), Labor Law and Labor Procedure (1,440),
Constitutional Law (1,439), Business Administration (1,305) and Logic (1,172).
The tail is long: 15 subjects have fewer than 100 questions and seven fewer
than 50, the smallest being Public Safety and Physical Security (13),
Human--Computer Interaction (23), Linguistics (37), Human Rights (37)
and Electoral Law (41). Only 17 subjects occur at high-school level.

\begin{table*}[p]
\centering
\scriptsize
\setlength{\tabcolsep}{3pt}
\begin{tabular}{@{}lrrr@{}}
\toprule
Subject & HS & UG & Total \\
\midrule
\multicolumn{4}{@{}l}{\emph{\textbf{Languages, Literature, Arts, and Communication}}} \\
Portuguese Language & 1,042 & 581 & 1,623 \\
English Language & 683 & 300 & 983 \\
Spanish Language & 6 & 200 & 206 \\
Literature & 475 & 93 & 568 \\
Linguistics & -- & 37 & 37 \\
Arts & 87 & 119 & 206 \\
Communication Studies & 27 & 361 & 388 \\
\midrule
\multicolumn{4}{@{}l}{\emph{\textbf{Mathematics, Statistics, and Logic}}} \\
Mathematics & 1,749 & 812 & 2,561 \\
Logic & 1,062 & 110 & 1,172 \\
Probability and Statistics & -- & 260 & 260 \\
\midrule
\multicolumn{4}{@{}l}{\emph{\textbf{Computing, Engineering, Architecture, and Design}}} \\
Programming, Algorithms and Data Structures & 80 & 522 & 602 \\
Computer Systems and Architecture & -- & 250 & 250 \\
Theory of Computation & -- & 110 & 110 \\
Software Engineering and Databases & -- & 610 & 610 \\
Computer Networks and Distributed Systems & -- & 379 & 379 \\
Artificial Intelligence, Data Science and Computer Vision & -- & 91 & 91 \\
Information Security & -- & 135 & 135 \\
Engineering & -- & 692 & 692 \\
Architecture and Urban Planning & -- & 192 & 192 \\
Design & -- & 218 & 218 \\
Computer Graphics & -- & 64 & 64 \\
Human--Computer Interaction & -- & 23 & 23 \\
\midrule
\multicolumn{4}{@{}l}{\emph{\textbf{Natural, Earth, Agricultural and Veterinary Sciences}}} \\
Physics & 879 & 184 & 1,063 \\
Chemistry & 566 & 250 & 816 \\
Biology & 568 & 159 & 727 \\
Earth and Environmental Sciences & -- & 234 & 234 \\
Agronomy & -- & 262 & 262 \\
Veterinary Medicine & -- & 119 & 119 \\
Animal Science & -- & 160 & 160 \\
\midrule
\multicolumn{4}{@{}l}{\emph{\textbf{Humanities}}} \\
History & 641 & 195 & 836 \\
Geography & 516 & 152 & 668 \\
Philosophy & 114 & 265 & 379 \\
Sociology & 193 & 204 & 397 \\
Psychology & -- & 270 & 270 \\
Education and Pedagogy & -- & 395 & 395 \\
Theology and Religious Studies & -- & 53 & 53 \\
\bottomrule
\end{tabular}
\hfill
\begin{tabular}{@{}lrrr@{}}
\toprule
Subject & HS & UG & Total \\
\midrule
\multicolumn{4}{@{}l}{\emph{\textbf{Social and Applied Social Sciences}}} \\
Political Science & -- & 43 & 43 \\
International Relations & -- & 146 & 146 \\
Public Policy and Public Administration & -- & 260 & 260 \\
Social Work & -- & 127 & 127 \\
Information Science & -- & 500 & 500 \\
Tourism & -- & 124 & 124 \\
Professional Ethics & -- & 914 & 914 \\
Gastronomy and Culinary Arts & -- & 97 & 97 \\
Public Safety and Physical Security & -- & 13 & 13 \\
\midrule
\multicolumn{4}{@{}l}{\emph{\textbf{Economics, Business, and Accounting}}} \\
Economics & -- & 539 & 539 \\
Business Administration & -- & 1,305 & 1,305 \\
Finance and Financial Systems & -- & 399 & 399 \\
Accounting & -- & 1,097 & 1,097 \\
\midrule
\multicolumn{4}{@{}l}{\emph{\textbf{Law}}} \\
Constitutional Law & -- & 1,439 & 1,439 \\
Administrative Law & -- & 949 & 949 \\
Civil Law and Civil Procedure & -- & 2,317 & 2,317 \\
Criminal Law and Criminal Procedure & -- & 1,749 & 1,749 \\
Labor Law and Labor Procedure & -- & 1,440 & 1,440 \\
Business Law & -- & 898 & 898 \\
Tax and Financial Law & -- & 798 & 798 \\
International Law & -- & 283 & 283 \\
Environmental Law & -- & 222 & 222 \\
Consumer Law & -- & 147 & 147 \\
Electoral Law & -- & 41 & 41 \\
Social Security Law & -- & 42 & 42 \\
Human Rights & -- & 37 & 37 \\
Child and Adolescent Law & -- & 113 & 113 \\
\midrule
\multicolumn{4}{@{}l}{\emph{\textbf{Health Sciences}}} \\
Internal Medicine & -- & 1,088 & 1,088 \\
General Surgery & -- & 503 & 503 \\
Pediatrics & -- & 415 & 415 \\
Obstetrics and Gynecology & -- & 365 & 365 \\
Family and Community Medicine and Public Health & -- & 469 & 469 \\
Biomedicine & -- & 64 & 64 \\
Nursing & -- & 98 & 98 \\
Pharmacy & -- & 102 & 102 \\
Physical Therapy and Speech-Language Pathology & -- & 217 & 217 \\
Nutrition & -- & 137 & 137 \\
Dentistry & -- & 113 & 113 \\
Physical Education & 17 & 146 & 163 \\
Occupational Therapy & -- & 53 & 53 \\
Radiology and Medical Imaging & -- & 167 & 167 \\
Aesthetics and Cosmetology & -- & 58 & 58 \\
\bottomrule
\end{tabular}
\caption{The 78 subjects of taxonomy 1.4 and the test questions assigned to each. \emph{HS}: high school; \emph{UG}: undergraduate; ``--'': no question of that level received the subject (the subject is not a candidate of any exam of that level).}
\label{tab:ka-short-subjects}
\end{table*}

\paragraph{By exam.} Figure~\ref{fig:ka-short-exams} shows the macro areas
of each exam in the test set. The exams with a declared scope stay inside it:
OAB is 91.5\% Law and 7.8\% Social and Applied Social Sciences
(Professional Ethics, a candidate of the bar exam), ENAM 97.1\% Law,
REVALIDA 98.2\% and the residency exams 94.0\% Health Sciences,
CFCES 88.6\% Accounting, OBI 93.9\% Mathematics,
Statistics and Logic (80.6\% Logic) and POSCOMP 69.7\% Computing and
30.3\% Mathematics, Statistics and Logic. The broad exams use their
lists widely: ENADE uses 77 of its 78 candidates and all nine areas, with
Business Administration first at only 9.7\%, BNDES uses 35 of 36
candidates, CNU 39 of 51 and BACEN all 26. The vestibular exams spread over
the school subjects, with no subject above 24\% in ENEM, FUVEST, BLUEX and
COMVEST, while the military academies split between languages and exact
sciences (AFA: Portuguese 26.5\%, Mathematics 25.6\%; IME:
English 32.2\%, Mathematics 28.4\%). Since the exam and the
edition are shown to the model, these shares describe the labels, not an
independent check of them.

\begin{figure*}[t]
\centering
\includegraphics[width=\textwidth]{\KnowledgeAreaFigDir/macro_area_by_exam.pdf}
\caption{Share of each exam's questions assigned to every macro area (\%) in the test set. Exams are grouped by academic level and sorted by size; empty cells are macro areas without any question of that exam, and ``$<$1'' marks shares below one percent.}
\label{fig:ka-short-exams}
\end{figure*}

\paragraph{Post-annotation.} The eight abstentions of the final run were
read by the authors. Six received a manual subject, always one of the exam's
candidates: in two cases the adjudicator had rejected the subject of one pass,
claiming that it was not in the list when it was, and that subject was
restored. In three cases the question's discipline is missing from the exam's
list by design (no law subject in ENEM or COMVEST, no administrative law in
the residency exams), and the closest candidate was chosen, and one is a boundary
case between General Surgery and Pediatrics. These rows carry the status
\emph{manual}, a null confidence, and the two model passes for audit. Two
questions were removed because their supporting chart or table is missing,
joining the 17 questions removed by the dataset revision for the same family
of defects. The annotated source corpus has 41,634 questions, 38,337 accepted, 3,291
adjudicated, and 6 manual, with no \texttt{UNCERTAIN}, and was published as
\textcolor{red}{\texttt{bench-temp-2/mmlu-pt-knowledge-annotated} at revision
\texttt{fb9d509a}\ldots}.
The published test set retains 34,855 accepted, 2,939 adjudicated and six
manual annotations, with no \texttt{UNCERTAIN}.

% -----------------------------------------------------------------------------
\subsection{Cost and reproducibility}
\label{app:ka-short-cost}

Table~\ref{tab:ka-short-cost} summarizes the final annotation run. The 41,636 questions
were annotated on one NVIDIA B200 in 03:38:39 (HH:mm:ss) from the start
of the job to the end of the export, of which 00:01:52 were spent loading the
weights and initializing the engine and 03:34:52 inside generation batches.
Each of the two passes sends 148.4 million prompt tokens, 3,564 per question
on average (2,020 at high-school and 4,039 at undergraduate level, where the
candidate lists are longer; the longest prompt has 10,367 tokens), and
receives about 72 completion tokens, the JSON object. Prefix caching was
disabled, so the prompt-token count is the full cost of re-reading the rules
and definitions for every question and an upper bound for a re-run with
caching. The run manifest records the source revision and its hash, the model
identifier and resolved revision, the taxonomy and alias hashes, the prompt
version, the code hashes, the library versions and the engine settings, and
the checkpoint stores the seed of every request. Sampling at temperature~1
with per-request seeds reproduces the same labels only on the same hardware,
kernels and batch composition; on another stack the labels should be expected
to differ at the level of the within-run disagreement.

\begin{table}[t]
\centering
\footnotesize
\setlength{\tabcolsep}{4pt}
\resizebox{\linewidth}{!}{
\begin{tabular}{@{}lrrrr@{}}
\toprule
Stage & Requests & Prompt tokens & Compl.\ tokens & Batch time\\
\midrule
Pass 1 & 41,636 & 148,407,132 & 2,981,987 & 01:43:19 \\
Pass 2 & 41,636 & 148,407,132 & 2,981,000 & 01:41:57 \\
Adjudication & 3,299 & 14,429,317 & 255,387 & 00:09:36 \\
Total & 86,571 & 311,243,581 & 6,218,374 & 03:34:52 \\
\bottomrule
\end{tabular}
}
\caption{Cost of the final run on one NVIDIA B200 (191.5~GB): FP8 weights of 115.37~GiB, vLLM 0.25.0 with XGrammar and CUDA graphs, batches of 128 questions, 50 concurrent sequences, context of 20,480 tokens, output budget of 8,192 tokens, prefix caching off. \emph{Requests}: successful model requests; the six requests retried after a schema violation are excluded from the counts and included in the batch time. \emph{Batch time}: sum of the durations of the generation batches; it excludes engine start-up and export.}
\label{tab:ka-short-cost}
\end{table}

% -----------------------------------------------------------------------------
\subsection{Limitations}
\label{app:ka-short-limitations}

Apart from the six manual assignments, there is no human-labelled reference
set against which to measure annotation accuracy. Agreement between passes,
stability between versions, and the shares per exam
measure the consistency of one model with itself, and the audits that
motivated the versions are engineering reviews of the model's output. The two passes disagree on about 8\% of the
questions, and the adjudicator is the same model, so a change between versions
is interpretable only where it exceeds that level, and versions 1.3 and 1.4
were tested on 23 and 15 questions before the third full run. The exam and
the edition are shown to the model, so the per-exam shares are not an
independent check of the labels. The reported confidence is a number the
model writes: 21 of the 23 abstentions of run 1.2 carried a confidence of 0.90
or more. Interdisciplinary questions receive one subject, and the candidate
lists of the vestibular and residency exams do not contain every discipline
their questions touch, as three of the six manual assignments show.
